% =====================================================
% CAPA / ABERTURA
% =====================================================
\capaRoteiro{\faCogs}{Monitoria de Bancada, Logística \& Moodle}
  {Guia operacional da Frente 4 para as Oficinas de 120 minutos da SNCT 2026}
  {\chip{\faLayerGroup}{Extensão A2 (20\%)}\hspace{4pt}
   \chip{\faTools}{Laboratório · Moodle IFSC · Certificação}\hspace{4pt}
   \chip{\faClipboardCheck}{Equipes de 3 estudantes}}

\vspace{0.4em}
\cartoesUC{Plano de bancada do laboratório, inscrição/acesso no Moodle, pesquisa de satisfação e evidências para o relatório}

\begin{boxObjetivos}
\begin{itemize}
  \item \textbf{Garantir a infraestrutura das bancadas do laboratório}, testando máquinas, rede, projetor, contas de contingência e pen-drives do Plano B antes do público entrar;
  \item \textbf{Configurar e validar o acesso dos participantes à sala do Moodle institucional} (autoinscrição com chave \texttt{SNCT2026} e acesso de visitante), atuando com papel de \textbf{Professor};
  \item \textbf{Criar a atividade de Pesquisa (Feedback) de 3 perguntas e a Lista de Presença} na Seção 4 do Moodle (com formulário impresso de reserva);
  \item \textbf{Coordenar a escala de monitoria durante os 120 minutos} (1 monitor para cada 4 a 5 participantes), aplicando a regra pedagógica de \emph{nunca tomar o teclado do aluno};
  \item \textbf{Consolidar as evidências da ação extensionista} (presenças, gráficos de satisfação do Moodle e registro fotográfico) para o relatório final.
\end{itemize}
\end{boxObjetivos}

\begin{mapaAula}
\etapaMapa{azulMedio}{1}{Missão desta semana}{Candidatura até 27/10 e divisão das frentes logísticas (Seção~\ref{sec:semana})}
\etapaMapa{verdeVivo}{2}{Acesso e Feedback no Moodle}{Chave de inscrição, contas de reserva e pesquisa de 3 perguntas (Seção~\ref{sec:moodle})}
\etapaMapa{ambar}{3}{Protocolo de Bancada}{Postura do monitor, escala 1:4 e roteiro minuto a minuto (Seção~\ref{sec:bancada})}
\etapaMapa{roxo}{4}{Marcos e evidências}{Lista de presença, termos de imagem e checklist final (Seção~\ref{sec:marcos})}
\end{mapaAula}

% =====================================================
% 1. ONDE ESTA FRENTE SE ENCAIXA E TAREFAS DA SEMANA
% =====================================================
\section{Onde esta frente se encaixa e o que fazer nesta semana}\label{sec:semana}

O melhor roteiro e o melhor código falham se, no início da oficina, metade dos computadores não abrir a sala do curso ou se um participante iniciante travar no primeiro erro de digitação e ficar sem ajuda. A \textbf{Frente 4 (Monitoria de Bancada \& Logística)} é a responsável por transformar o laboratório em um ambiente acolhedor, fluido e sem gargalos técnicos.

\begin{table}[H]
\centering
\small
\renewcommand{\arraystretch}{1.45}\rowcolors{2}{fundo}{white}
\begin{tabularx}{\textwidth}{RYY}
\cabecalho \cab{Frente de Trabalho} & \cab{Entregável Principal} & \cab{Conexão com a Frente de Monitoria} \\
1. Didática \& Micro-Roteiro & PDF de até 4 páginas + Sala no Moodle IFSC. & Vocês distribuem os roteiros impressos nas bancadas e testam a sala do Moodle. \\
2. Starter Kit \& GitHub & Repositório template + \texttt{starter-kit.zip}. & Vocês recebem os pen-drives do Plano B e deixam o kit pré-testado nos PCs do lab. \\
3. Slides-Guia (8 Telas) & Apresentação HTML no GitHub Pages. & Vocês enviam o QR Code da Pesquisa de Satisfação para entrar na Tela 8 dos slides. \\
\textbf{4. Monitoria \& Logística (Você)} & Preparação do lab, inscrição no Moodle e avaliação. & Garante atendimento nas ilhas, presença, satisfação e fotos para a extensão. \\
\end{tabularx}
\caption{Articulação entre a Frente 4 e as demais equipes da oficina.}
\end{table}

\subsection*{Pipeline Visual: Da Produção em Equipe à Aplicação nos Dias da SNCT}

O diagrama abaixo resume o fluxo completo de execução que conecta o trabalho desta semana ao dia da oficina presencial no laboratório:

\begin{center}
\begin{adjustbox}{max width=\linewidth}
\begin{tikzpicture}[
    node distance=0.45cm,
    etapa/.style={rectangle, rounded corners=6pt, draw=#1, fill=#1!8, line width=1pt,
                  text width=2.85cm, minimum height=2.35cm, align=center, inner sep=5pt, font=\scriptsize\sffamily},
    seta/.style={-{Latex[length=2.6mm, width=2mm]}, line width=1.2pt, color=azul}
]
\node[etapa=azulMedio] (f1) {
  {\bfseries\color{azulMedio}1. PRODUÇÃO}\\[2pt]
  \textbf{Frentes 1 a 4}\\
  Codificação do Starter Kit, Slides (8 telas), Roteiro A4 e Pesquisa.\\[2pt]
  {\tiny\color{suave}(Semana 28--30/10)}
};

\node[etapa=ambar, right=of f1] (f2) {
  {\bfseries\color{ambar!85!black}2. TESTE NO LAB}\\[2pt]
  \textbf{Revisão Cruzada}\\
  Teste real nas máquinas do lab, código 100\% igual ao PDF e Plano B.\\[2pt]
  {\tiny\color{suave}(Pré-publicação)}
};

\node[etapa=verdeVivo, right=of f2] (f3) {
  {\bfseries\color{verde}3. MOODLE IFSC}\\[2pt]
  \textbf{Alunos = Professores}\\
  Criação das 5 seções nas Salas do Curso 1 e do Curso 2 no Moodle.\\[2pt]
  {\tiny\color{suave}(Chave: SNCT2026)}
};

\node[etapa=ciano, right=of f3] (f4) {
  {\bfseries\color{ciano}4. KIT + CÓDIGOS}\\[2pt]
  \textbf{Disponibilização}\\
  Upload do \texttt{.zip}, PDF, link GitHub Template e Slides na sala.\\[2pt]
  {\tiny\color{suave}(Pronto p/ público)}
};

\node[etapa=roxo, right=of f4] (f5) {
  {\bfseries\color{roxo}5. DIAS DA SNCT}\\[2pt]
  \textbf{Aplicação (120 min)}\\
  Oficina nas bancadas, site no celular (QR Code) e avaliação final.\\[2pt]
  {\tiny\color{suave}(Evento + Relatório)}
};

\draw[seta] (f1.east) -- (f2.west);
\draw[seta] (f2.east) -- (f3.west);
\draw[seta] (f3.east) -- (f4.west);
\draw[seta] (f4.east) -- (f5.west);
\end{tikzpicture}
\end{adjustbox}
\end{center}


\begin{atencao}{Ação imediata para esta semana (ausência docente em sala)}
\begin{enumerate}[leftmargin=1.4em, itemsep=0.25em]
  \item \textbf{Até amanhã (terça-feira, 27/10, às 23h59):} acesse o painel da disciplina com seu login Google e candidate-se na frente \emph{Monitoria de Bancada \& Logística} do Curso 1 ou do Curso 2.
  \item \textbf{Durante a aula desta semana:} dividam os três papéis abaixo, configurem a Pesquisa de Satisfação na sala do Moodle (ou Google Forms espelho) e registrem o link no painel.
\end{enumerate}
\end{atencao}

\begin{table}[H]
\centering
\small
\renewcommand{\arraystretch}{1.45}\rowcolors{2}{fundo}{white}
\begin{tabularx}{\textwidth}{RYY}
\cabecalho \cab{Papel na Equipe} & \cab{Responsabilidade Técnica} & \cab{Entrega até o fim desta semana} \\
Membro A — Acessos \& Moodle IFSC & Configurar autoinscrição/visitante no Moodle, criar a Pesquisa de Satisfação e testar login nas bancadas. & Pesquisa criada na Seção 4 do Moodle e chave \texttt{SNCT2026} testada. \\
Membro B — Infraestrutura \& Plano B & Auditar os computadores do laboratório (VS Code, navegador, rede) e preparar contas GitHub de reserva. & Mapa das máquinas do laboratório e 3 contas de contingência prontas. \\
Membro C — Credenciamento \& Relatório & Preparar a lista de presença, os termos de autorização de imagem e o roteiro de fotos da oficina. & Modelos de presença e autorização prontos para impressão e no Moodle. \\
\end{tabularx}
\caption{Divisão de trabalho recomendada para os 3 estudantes da Frente 4.}
\end{table}

% =====================================================
% 2. INSCRIÇÃO E AVALIAÇÃO NO MOODLE INSTITUCIONAL
% =====================================================
\section{Inscrição, Acesso de Visitantes e Avaliação no Moodle}\label{sec:moodle}

Como os estudantes da turma possuem permissão de \textbf{Professor} no Moodle institucional do IFSC, a Frente 4 deve garantir que nenhum participante perca tempo na tela de login.

\begin{dica}{Ações separadas por curso: Monitoria do Curso 1 vs. Monitoria do Curso 2}
Cada oficina possui a sua própria sala no Moodle institucional e o seu próprio trio de monitoria:
\begin{itemize}[leftmargin=1.4em, itemsep=0.2em]
  \item \textbf{Equipe de Monitoria do Curso 1 (Meu Primeiro Site no Ar):} prepara e testa o pacote \texttt{starter-kit-snct.zip}, valida o link \texttt{wa.me} e aplica a lista de presença/avaliação da Sala do Curso 1.
  \item \textbf{Equipe de Monitoria do Curso 2 (Painel Verde: Coleta Seletiva \& Ecopontos):} prepara e testa o pacote \texttt{starter-kit-painel-verde.zip}, confere nas bancadas o fechamento de chaves/vírgulas do array \texttt{LOCAIS\_COLETA} em \texttt{js/app.js} e aplica a avaliação da Sala do Curso 2.
\end{itemize}
\end{dica}

\passo{2.1}{Habilitar duplo caminho de acesso na sala do seu curso no Moodle}
\begin{enumerate}[leftmargin=1.4em, itemsep=0.25em]
  \item \textbf{Para estudantes do IFSC (com login institucional):} em \botao{Participantes}\seta\botao{Métodos de inscrição}, ative \textbf{Autoinscrição (Estudante)} com a chave de inscrição \texttt{SNCT2026};
  \item \textbf{Para público externo da comunidade (sem conta no Moodle):} ative o método \textbf{Acesso de visitante} na sala do curso para que qualquer pessoa consiga baixar o PDF e o \texttt{starter-kit-snct.zip} na hora, e deixe um link espelho no Google Forms para o questionário final.
\end{enumerate}

\manterJunto{26}
\passo{2.2}{Criar a Pesquisa de Satisfação de 3 perguntas na Seção 4 do Moodle}

Na \textbf{Seção 4 — Avaliação \& Certificado}, adicione a atividade \botao{Pesquisa} (ou \emph{Feedback}) com exatamente 3 perguntas rápidas para que 100\% da turma responda nos últimos 10 minutos:

\begin{codigo}[txt]{Perguntas da Pesquisa de Satisfacao (Seção 4 do Moodle / Google Forms)}
PERGUNTA 1 (Escala de 1 a 5):
De 1 a 5, o quanto o Micro-Roteiro e a explicacao ajudaram voce a colocar seu projeto no ar?
[ ] 1   [ ] 2   [ ] 3   [ ] 4   [ ] 5 (Excelente)

PERGUNTA 2 (Multipla escolha):
Voce conseguiu abrir o projeto publicado no seu smartphone via QR Code?
[ ] Sim, abriu perfeitamente no celular!
[ ] Abriu no computador do laboratorio, mas nao testei no celular.
[ ] Tive dificuldade na etapa do GitHub Pages.

PERGUNTA 3 (Resposta curta opcional):
O que voce mais gostou na oficina de hoje e que tema gostaria de ver na proxima edicao?
\end{codigo}

% =====================================================
% 3. PROTOCOLO DE MONITORIA DE BANCADA (120 MINUTOS)
% =====================================================
\section{Protocolo de Monitoria de Bancada e Escala do Dia}\label{sec:bancada}

\begin{conceito}{Regra de Ouro da Monitoria: Mãos Atrás das Costas!}
Quando um participante chamar ajuda na bancada, \textbf{nunca pegue o mouse nem digite no teclado por ele}. Aponte na tela onde está o detalhe, mostre a linha correspondente no Micro-Roteiro impresso e deixe que a própria pessoa faça a alteração. Quando o próprio participante resolve e vê a tela atualizar, a aprendizagem realmente acontece.
\end{conceito}

\begin{table}[H]
\centering
\small
\renewcommand{\arraystretch}{1.45}\rowcolors{2}{fundo}{white}
\begin{tabularx}{\textwidth}{RYY}
\cabecalho \cab{Janela de Tempo} & \cab{Ação no Telão (Instrutor)} & \cab{Ação da Monitoria nas Bancadas} \\
$-30$ a 00 min (Preparação) & Testar projetor, abrir Tela 1 dos slides e testar som/foco. & Ligar todos os PCs, abrir a sala do Moodle e colocar 1 roteiro impresso por máquina. \\
00--15 min (Acolhida) & Boas-vindas e demonstração do resultado final via QR Code. & Receber quem chega atrasado, recolher assinatura na lista de presença e termo de imagem. \\
15--35 min (Bloco 1) & Guiar o download do \texttt{starter-kit.zip} e edição do \texttt{:root}. & Destravar extração do \texttt{.zip} e abertura da pasta no VS Code (usar pen-drive se a rede lentificar). \\
35--70 min (Bloco 2) & Explicar a customização central (HTML/WhatsApp ou JS). & Circular pelas ilhas conferindo fechamento de aspas e tags antes do Bloco 3. \\
70--110 min (Blocos 3 e 4) & Demonstrar upload no GitHub, Pages e geração do QR Code. & Apoiar login no GitHub (fornecer conta de contingência a quem não receber código 2FA). \\
110--120 min (Fechamento) & Exibir Tela 8 com QR Code da Pesquisa de Satisfação. & Garantir que todos enviem a avaliação no Moodle/Forms e registrar a foto oficial da turma. \\
\end{tabularx}
\caption{Roteiro operacional minuto a minuto para o ensaio geral e para o dia da SNCT.}
\end{table}

% =====================================================
% 4. CRONOGRAMA DE MARCOS E CHECKLIST DE ENTREGA
% =====================================================
\section{Cronograma de Marcos e Checklist de Entrega}\label{sec:marcos}

\begin{checklist}[Checklist de prontidão da Frente 4 (Monitoria, Logística \& Moodle)]
\begin{itemize}[label={\caixaMarcar}, itemsep=0.35em]
  \item Os 3 estudantes registraram candidatura na frente até 27/10 às 23h59 no painel da disciplina;
  \item A autoinscrição com chave \texttt{SNCT2026} e o acesso de visitante estão ativos na sala do Moodle;
  \item A Pesquisa de Satisfação de 3 perguntas está publicada na Seção 4 do Moodle (e o link foi enviado à Frente 3);
  \item A lista de presença impressa e os termos de autorização de uso de imagem estão diagramados e impressos;
  \item Pelo menos 3 contas de contingência do GitHub e 2 pen-drives com o Plano B completo foram testados no laboratório;
  \item A escala de quem fica no telão e quem cuida de cada fila de bancadas foi definida para o ensaio e para a SNCT;
  \item O link da pasta de logística/Moodle foi registrado na Central de Demandas da disciplina.
\end{itemize}
\end{checklist}

\vfill
{\centering\footnotesize\color{suave}\faMapMarkerAlt\hspace{0.4em}IFSC Câmpus Garopaba — Programação para a Internet 2 (Curricularização da Extensão SNCT 2026)\par}

